\subsection{GRADED ALGEBRA ASSOCIATED WITH A FILTERED ALGEBRA}

Let G be a commutative group with a real filtration \((G_{\alpha})_{\alpha \in \mathbb{R}}\) . As before we write

\[
\mathrm{G} _ {\alpha} ^ {+} = \bigcup_ {\beta > \alpha} \mathrm{G} _ {\beta};\tag{9}
\]

clearly \(\mathrm{G}_{\alpha}^{+}\) is a subgroup of \(\mathrm{G}_{\alpha}\) . We write \(\mathrm{gr}_{\alpha}(\mathrm{G}) = \mathrm{G}_{\alpha} / \mathrm{G}_{\alpha}^{+}\) and

\[
\operatorname{gr} (\mathrm{G}) = \bigoplus_ {\alpha \in \mathbf {R}} \operatorname{gr} _ {\alpha} (\mathrm{G}).
\]

The graded group associated with the filtered group G is the group \(\operatorname{gr}(G)\) with its natural graduation of type R.

Remark. When the filtration \((\mathrm{G}_{\alpha})\) is integral, \(\mathrm{gr}_{\alpha}(\mathrm{G}) = \{0\}\) for non-integral \(\alpha\) and \(\mathrm{gr}_n(\mathrm{G}) = \mathrm{G}_n / \mathrm{G}_{n - 1}\) for every integer \(n\) . The definition of the associated graded group therefore coincides essentially with that of Commutative Algebra, Chapter III, § 2, no. 3.

Let \(\mathbf{A}\) be an algebra (resp. unital algebra) and \((\mathbf{A}_{\alpha})_{\alpha \in \mathbb{R}}\) a real filtration compatible with the algebra (resp. unital algebra) structure (no. 1). Then

\[
\mathrm{A} _ {\alpha}. \mathrm{A} _ {\beta} \subset \mathrm{A} _ {\alpha + \beta}, \quad \mathrm{A} _ {\alpha} ^ {+}. \mathrm{A} _ {\beta} + \mathrm{A} _ {\alpha}. \mathrm{A} _ {\beta} ^ {+} \subset \mathrm{A} _ {\alpha + \beta} ^ {+},
\]

and the bilinear mapping of \(A_{\alpha} \times A_{\beta}\) into \(A_{\alpha+\beta}\) the restriction of multiplication on A defines on taking quotients a bilinear mapping

\[
\operatorname{gr} _ {\alpha} (\mathrm{A}) \times \operatorname{gr} _ {\beta} (\mathrm{A}) \rightarrow \operatorname{gr} _ {\alpha + \beta} (\mathrm{A}).
\]

We derive a bilinear mapping of \(\mathrm{gr}(\mathbf{A}) \times \mathrm{gr}(\mathbf{A})\) into \(\mathrm{gr}(\mathbf{A})\) which makes it a graded algebra (resp. unital graded algebra) of type R. If A is an associative (resp. commutative, resp. Lie) algebra, so is \(\mathrm{gr}(\mathbf{A})\) .

